\section{Kết luận}

\begin{frame}
\frametitle{Tổng kết phương pháp}
\framesubtitle{Vehicle detection, tracking and counting}

\begin{itemize}[<+->]
  \item Đề xuất chuỗi xử lý gồm ba giai đoạn: phát hiện, theo dõi và đếm
        phương tiện
  \item Sử dụng YOLOv8n làm bộ phát hiện phương tiện trong từng khung hình
  \item Sử dụng BYTE để duy trì định danh phương tiện qua thời gian
  \item Đếm phương tiện dựa trên ID theo dõi và ROI do người dùng cấu hình
  \item Xây dựng pipeline và công cụ minh họa có thể chạy từ dòng lệnh
\end{itemize}

\end{frame}

\begin{frame}
\frametitle{Kết quả thực nghiệm}
\framesubtitle{Detection và tracking}

\begin{itemize}[<+->]
  \item YOLOv8n sau fine-tuning hội tụ ổn định trên UA-DETRAC
  \item Kết quả detection trên tập kiểm tra đạt Precision $0.7352$,
        Recall $0.5488$ và mAP@50 $0.5003$
  \item BYTE cao hơn SORT trên tập xác thực ở các độ đo chính như HOTA, MOTA
        và IDF1
  \item YOLOv8n + BYTE trên tập kiểm tra đạt HOTA $60.716$, MOTA $64.676$ và
        IDF1 $70.345$
  \item Các kết quả cho thấy hướng tracking-by-detection là khả thi cho video
        giao thông
\end{itemize}

\end{frame}

\begin{frame}
\frametitle{Những điều chứng minh được}
\framesubtitle{Ý nghĩa của kết quả}

\begin{itemize}[<+->]
  \item YOLOv8n học được đặc trưng phương tiện trên dữ liệu giao thông thực tế
  \item BYTE giúp duy trì ID tốt hơn SORT nhờ tận dụng cả phát hiện có độ tin
        cậy thấp
  \item Chất lượng tracking phụ thuộc mạnh vào chất lượng hộp phát hiện từ
        detector
  \item Phân tích HOTA cho thấy DetA giảm nhanh khi ngưỡng đánh giá nghiêm
        ngặt hơn
  \item Cải thiện detector là hướng quan trọng để nâng hiệu năng toàn hệ thống
\end{itemize}

\end{frame}

\begin{frame}
\frametitle{Ưu điểm và hạn chế}
\framesubtitle{Đánh giá phương pháp}

\begin{table}
  \centering
  \begin{tblr}{
    colspec = {X[1,l] X[1,l]},
    vline{2}, hline{2},
  }
    Ưu điểm & Hạn chế \\
    Cấu trúc đơn giản, dễ triển khai &
    Phụ thuộc vào chất lượng detection \\
    YOLOv8n nhỏ, tốc độ suy luận tốt &
    Hiệu năng chưa đều giữa các lớp \\
    BYTE giảm mất dấu và giữ ID tốt hơn SORT &
    Đối tượng nhỏ hoặc che khuất vẫn khó xử lý \\
    Dễ thay detector hoặc tracker &
    Counting cần đánh giá định lượng sâu hơn \\
  \end{tblr}
\end{table}

\end{frame}

\begin{frame}
\frametitle{Hướng phát triển}
\framesubtitle{Các cải tiến tiếp theo}

\begin{itemize}[<+->]
  \item Xử lý mất cân bằng dữ liệu bằng tăng cường dữ liệu hoặc lấy mẫu lại
  \item Thử nghiệm các biến thể YOLO lớn hơn hoặc tối ưu suy luận tốt hơn
  \item Bổ sung đặc trưng ngoại quan để giảm nhầm lẫn ID khi che khuất
  \item Đánh giá trực tiếp độ chính xác counting trên các ROI hoặc đường đếm
        có nhãn chuẩn
  \item Mở rộng hệ thống cho nhiều góc camera và các kịch bản giao thông phức
        tạp hơn
\end{itemize}

\end{frame}
